% GENERATED FILE. Edit canonical lessons, then run npm run generate:books.
% canonical-source-hash: fnv1a64:c4ad58fd5cb4837d
% canonical-lessons: KA-C524-cikkappa, KA-C524-doddappa, KA-C524-cikkamma, KA-C524-doddamma, KA-C524-sodaraliya

\chapter{Uncles and Aunts}
\label{ch:ka-uncles-and-aunts}
\hlchaptermodality{524}

\begin{chapteropening}
\textbf{By the end of this chapter:} \emph{I can say father's younger brother, father's elder brother, mother's younger sister, mother's elder sister and a nephew.}

Complete the last lesson of chapter 524: I can say father's younger brother, father's elder brother, mother's younger sister, mother's elder sister and a nephew.
\end{chapteropening}

\section[\texorpdfstring{\kn{ಚಿಕ್ಕಪ್ಪ} (cikkappa)}{ಚಿಕ್ಕಪ್ಪ (cikkappa)}]{\kn{ಚಿಕ್ಕಪ್ಪ} (cikkappa) — father's younger brother}

\label{lesson:KA-C524-cikkappa}



\begin{quote}
Before the new one: say the Kannada for the pain of parting, then the Kannada for married life.
\end{quote}

\subsection*{You'll want to know: \kn{ಚಿಕ್ಕಪ್ಪ}}
\textbf{\kn{ಚಿಕ್ಕಪ್ಪ}} — \emph{cikkappa} — \textquotedblleft{}father's younger brother\textquotedblright{}.

\begin{quote}
\kn{ನನ್ನ} \kn{ಚಿಕ್ಕಪ್ಪ} \kn{ಹಳ್ಳಿಯಲ್ಲಿ} \kn{ರೈತರು}.
\end{quote}

\begin{quote}
\emph{nanna cikkappa haḷḷiyalli raitaru.}
\end{quote}

\begin{quote}
My father's younger brother is a farmer in the village.
\end{quote}

\kn{ಚಿಕ್ಕ}, small + \kn{ಅಪ್ಪ}, father: a \textquotedblleft{}small father\textquotedblright{}.

\begin{cousinweb}
\textbf{\kn{ಚಿಕ್ಕಪ್ಪ}} (\emph{cikkappa}) is a native Dravidian compound: \emph{cikka}, \textquotedblleft{}small, young\textquotedblright{} (DEDR 2495), and \emph{appa}, \textquotedblleft{}father\textquotedblright{} (DEDR 156).
\end{cousinweb}

\begin{grammarlens}[title={What you've built}]
One more word for describing what you lived and felt.
\end{grammarlens}

\subsection*{Your turn}
\begin{itemize}
\raggedright
  \item \emph{Say it:} \emph{cikkappa}
  \item \emph{Say it:} \emph{cikkappa}, once more
  \item \emph{Say it:} \emph{nanna cikkappa haḷḷiyalli raitaru.} — My father's younger brother is a farmer in the village.
  \item \emph{Recall:} say the Kannada for the pain of parting, then the Kannada for married life, then say \emph{cikkappa} again
\end{itemize}

\begin{tcolorbox}[breakable,colback=teal!4,colframe=teal!35!black,title={Before you move on}]
What does \kn{ಚಿಕ್ಕಪ್ಪ} mean? (\textquotedblleft{}Father's younger brother\textquotedblright{}.) Say it once more.
\end{tcolorbox}
\section[\texorpdfstring{\kn{ದೊಡ್ಡಪ್ಪ} (doḍḍappa)}{ದೊಡ್ಡಪ್ಪ (doḍḍappa)}]{\kn{ದೊಡ್ಡಪ್ಪ} (doḍḍappa) — father's elder brother}

\label{lesson:KA-C524-doddappa}



\begin{quote}
Before the new one: say the Kannada for married life, then the Kannada for father's younger brother.
\end{quote}

\subsection*{You'll want to know: \kn{ದೊಡ್ಡಪ್ಪ}}
\textbf{\kn{ದೊಡ್ಡಪ್ಪ}} — \emph{doḍḍappa} — \textquotedblleft{}father's elder brother\textquotedblright{}.

\begin{quote}
\kn{ದೊಡ್ಡಪ್ಪ} \kn{ನಮ್ಮ} \kn{ಕುಟುಂಬದ} \kn{ಹಿರಿಯರು}.
\end{quote}

\begin{quote}
\emph{doḍḍappa namma kuṭumbada hiriyaru.}
\end{quote}

\begin{quote}
My father's elder brother is an elder of our family.
\end{quote}

\kn{ದೊಡ್ಡ}, big + \kn{ಅಪ್ಪ}, father: a \textquotedblleft{}big father\textquotedblright{}.

\begin{cousinweb}
\textbf{\kn{ದೊಡ್ಡಪ್ಪ}} (\emph{doḍḍappa}) is a native Dravidian compound: \emph{doḍḍa}, \textquotedblleft{}big, great\textquotedblright{} (DEDR 3491), and \emph{appa}, \textquotedblleft{}father\textquotedblright{} (DEDR 156).
\end{cousinweb}

\begin{grammarlens}[title={What you've built}]
One more word for describing what you lived and felt.
\end{grammarlens}

\subsection*{Your turn}
\begin{itemize}
\raggedright
  \item \emph{Say it:} \emph{doḍḍappa}
  \item \emph{Say it:} \emph{doḍḍappa}, once more
  \item \emph{Say it:} \emph{doḍḍappa namma kuṭumbada hiriyaru.} — My father's elder brother is an elder of our family.
  \item \emph{Recall:} say the Kannada for married life, then the Kannada for father's younger brother, then say \emph{doḍḍappa} again
\end{itemize}

\begin{tcolorbox}[breakable,colback=teal!4,colframe=teal!35!black,title={Before you move on}]
What does \kn{ದೊಡ್ಡಪ್ಪ} mean? (\textquotedblleft{}Father's elder brother\textquotedblright{}.) Say it once more.
\end{tcolorbox}
\section[\texorpdfstring{\kn{ಚಿಕ್ಕಮ್ಮ} (cikkamma)}{ಚಿಕ್ಕಮ್ಮ (cikkamma)}]{\kn{ಚಿಕ್ಕಮ್ಮ} (cikkamma) — mother's younger sister}

\label{lesson:KA-C524-cikkamma}



\begin{quote}
Before the new one: say the Kannada for father's younger brother, then the Kannada for father's elder brother.
\end{quote}

\subsection*{You'll want to know: \kn{ಚಿಕ್ಕಮ್ಮ}}
\textbf{\kn{ಚಿಕ್ಕಮ್ಮ}} — \emph{cikkamma} — \textquotedblleft{}mother's younger sister\textquotedblright{}.

\begin{quote}
\kn{ರಜೆಯಲ್ಲಿ} \kn{ನಾನು} \kn{ಚಿಕ್ಕಮ್ಮನ} \kn{ಮನೆಗೆ} \kn{ಹೋಗುತ್ತೇನೆ}.
\end{quote}

\begin{quote}
\emph{rajeyalli nānu cikkammana manege hōguttēne.}
\end{quote}

\begin{quote}
In the holidays I go to my mother's younger sister's house.
\end{quote}

Also the wife of \kn{ಚಿಕ್ಕಪ್ಪ}.

\begin{cousinweb}
\textbf{\kn{ಚಿಕ್ಕಮ್ಮ}} (\emph{cikkamma}) is a native Dravidian compound: \emph{cikka}, \textquotedblleft{}small, young\textquotedblright{} (DEDR 2495), and \emph{amma}, \textquotedblleft{}mother\textquotedblright{} (DEDR 183).
\end{cousinweb}

\begin{grammarlens}[title={What you've built}]
One more word for describing what you lived and felt.
\end{grammarlens}

\subsection*{Your turn}
\begin{itemize}
\raggedright
  \item \emph{Say it:} \emph{cikkamma}
  \item \emph{Say it:} \emph{cikkamma}, once more
  \item \emph{Say it:} \emph{rajeyalli nānu cikkammana manege hōguttēne.} — In the holidays I go to my mother's younger sister's house.
  \item \emph{Recall:} say the Kannada for father's younger brother, then the Kannada for father's elder brother, then say \emph{cikkamma} again
\end{itemize}

\begin{tcolorbox}[breakable,colback=teal!4,colframe=teal!35!black,title={Before you move on}]
What does \kn{ಚಿಕ್ಕಮ್ಮ} mean? (\textquotedblleft{}Mother's younger sister\textquotedblright{}.) Say it once more.
\end{tcolorbox}
\section[\texorpdfstring{\kn{ದೊಡ್ಡಮ್ಮ} (doḍḍamma)}{ದೊಡ್ಡಮ್ಮ (doḍḍamma)}]{\kn{ದೊಡ್ಡಮ್ಮ} (doḍḍamma) — mother's elder sister}

\label{lesson:KA-C524-doddamma}



\begin{quote}
Before the new one: say the Kannada for father's elder brother, then the Kannada for mother's younger sister.
\end{quote}

\subsection*{You'll want to know: \kn{ದೊಡ್ಡಮ್ಮ}}
\textbf{\kn{ದೊಡ್ಡಮ್ಮ}} — \emph{doḍḍamma} — \textquotedblleft{}mother's elder sister\textquotedblright{}.

\begin{quote}
\kn{ದೊಡ್ಡಮ್ಮ} \kn{ರುಚಿಯಾದ} \kn{ಹೋಳಿಗೆ} \kn{ಮಾಡುತ್ತಾರೆ}.
\end{quote}

\begin{quote}
\emph{doḍḍamma ruciyāda hōḷige māḍuttāre.}
\end{quote}

\begin{quote}
My mother's elder sister makes tasty holige.
\end{quote}

Also the wife of \kn{ದೊಡ್ಡಪ್ಪ}.

\begin{cousinweb}
\textbf{\kn{ದೊಡ್ಡಮ್ಮ}} (\emph{doḍḍamma}) is a native Dravidian compound: \emph{doḍḍa}, \textquotedblleft{}big\textquotedblright{} (DEDR 3491), and \emph{amma}, \textquotedblleft{}mother\textquotedblright{} (DEDR 183).
\end{cousinweb}

\begin{grammarlens}[title={What you've built}]
One more word for describing what you lived and felt.
\end{grammarlens}

\subsection*{Your turn}
\begin{itemize}
\raggedright
  \item \emph{Say it:} \emph{doḍḍamma}
  \item \emph{Say it:} \emph{doḍḍamma}, once more
  \item \emph{Say it:} \emph{doḍḍamma ruciyāda hōḷige māḍuttāre.} — My mother's elder sister makes tasty holige.
  \item \emph{Recall:} say the Kannada for father's elder brother, then the Kannada for mother's younger sister, then say \emph{doḍḍamma} again
\end{itemize}

\begin{tcolorbox}[breakable,colback=teal!4,colframe=teal!35!black,title={Before you move on}]
What does \kn{ದೊಡ್ಡಮ್ಮ} mean? (\textquotedblleft{}Mother's elder sister\textquotedblright{}.) Say it once more.
\end{tcolorbox}
\section[\texorpdfstring{\kn{ಸೋದರಳಿಯ} (sōdaraḷiya)}{ಸೋದರಳಿಯ (sōdaraḷiya)}]{\kn{ಸೋದರಳಿಯ} (sōdaraḷiya) — a nephew}

\label{lesson:KA-C524-sodaraliya}



\begin{quote}
Before the new one: say the Kannada for mother's younger sister, then the Kannada for mother's elder sister.
\end{quote}

\subsection*{You'll want to know: \kn{ಸೋದರಳಿಯ}}
\textbf{\kn{ಸೋದರಳಿಯ}} — \emph{sōdaraḷiya} — \textquotedblleft{}a nephew\textquotedblright{}.

\begin{quote}
\kn{ನನ್ನ} \kn{ಸೋದರಳಿಯ} \kn{ಕಾಲೇಜಿನಲ್ಲಿ} \kn{ಓದುತ್ತಾನೆ}.
\end{quote}

\begin{quote}
\emph{nanna sōdaraḷiya kālējinalli ōduttāne.}
\end{quote}

\begin{quote}
My nephew studies at college.
\end{quote}

A sister's son: \kn{ಸೋದರ}, sibling + \kn{ಅಳಿಯ}. A niece is \kn{ಸೋದರ} \kn{ಸೊಸೆ}.

\begin{cousinweb}
\textbf{\kn{ಸೋದರಳಿಯ}} (\emph{sōdaraḷiya}) joins Sanskrit \emph{sodara}, \textquotedblleft{}born of the same womb\textquotedblright{} (Monier-Williams), and the native \emph{aḷiya}, \textquotedblleft{}son-in-law; sister's son\textquotedblright{} (DEDR 2410).
\end{cousinweb}

\begin{grammarlens}[title={What you've built}]
One more word for describing what you lived and felt.
\end{grammarlens}

\subsection*{Your turn}
\begin{itemize}
\raggedright
  \item \emph{Say it:} \emph{sōdaraḷiya}
  \item \emph{Say it:} \emph{sōdaraḷiya}, once more
  \item \emph{Say it:} \emph{nanna sōdaraḷiya kālējinalli ōduttāne.} — My nephew studies at college.
  \item \emph{Recall:} say the Kannada for mother's younger sister, then the Kannada for mother's elder sister, then say \emph{sōdaraḷiya} again
\end{itemize}

\begin{tcolorbox}[breakable,colback=teal!4,colframe=teal!35!black,title={Before you move on}]
What does \kn{ಸೋದರಳಿಯ} mean? (\textquotedblleft{}A nephew\textquotedblright{}.) Say it once more.
\end{tcolorbox}
